\section*{Working Scope and Reading Route}
\addcontentsline{toc}{section}{Working Scope and Reading Route}
\label{v480:reading-route}

This is the full-source version v\workingversion{}, derived from the
complete v4.7.0 monolith, not from the older modular synchronization
baseline, and continuing the published v4.8.0 line, whose layers are
carried in (\S\ref{v480:support-release-layers} and
\S\ref{v480:support-post-release-layers}). Every byte of the inherited
source is retained in uniquely included fragments.
Provenance of the previous versions: v4.8.0, Version DOI
\href{https://doi.org/10.5281/zenodo.23365244}{10.5281/zenodo.23365244};
v4.7.0, Version DOI
\href{https://doi.org/10.5281/zenodo.20232911}{10.5281/zenodo.20232911}.
The Concept DOI on the title page is common to all versions. The main line brings together the core
construction, physical readouts, explicit premise qualifications and
the next research interface. The appendices retain the complete
definitions, proofs, protocols, finite ledgers and provenance.

\paragraph{What changes, and what does not.}
The reading route changes; a stored or relocated assertion does not
thereby become a valid theorem. The inherited material is preserved
as source evidence, including unresolved or overstrong historical
wording. The explicit qualifications in
Section~\ref{v480:premise-audit} govern the known MM/HC, exponent and
OP5 issues. Section~\ref{v480:koide-kernel} proves the additional Koide condition
and constructs a relative local kernel model, with its remaining native
activation obligations made explicit. Section~\ref{v480:hc-mass-identification}
tests the HC-to-lepton-root identification, distinguishes sector and
generation triples, and isolates the missing residual coupling.
Section~\ref{v480:generation-coupling} constructs a conditional
closure-defect coupling and closed generation seeds, separates slot
readouts from candidate eigenmasses, and tests the pacer phase.
Section~\ref{v480:native-operator-link} tests the native operator
bridge, derives the minimal penalty compression, and separates a
conditional gapped-sector mass law from exact closure.
Section~\ref{v480:coupling-scale-candidate} evaluates a normalized
cyclic coupling candidate, tests the scale choices and proves a
fixed-shape mass-ratio obstruction.
Section~\ref{v480:tensor-shift-coupling} extends this test to explicit
comparison and kinetic metrics, distinguishing covariant shifts
from added physical tensor weights and target-derived fits.
Section~\ref{v480:hc-dynamics-geometry} carries the conditional HC
program through projected dynamics, real actions and bilinear geometry.
It separates the off-shell null stress match from the current test and
the conditional same-induced-frame trace obstruction.
Section~\ref{v490:cartan-source} adds the explicitly declared neutral
Cartan comparison, conditional gravitational constraint propagation,
and a finite normal-source law for compatible tangential jets. These
results do not establish the native action or regional existence.
The current endpoint is summarized in
Section~\ref{v480:current-v4}; the v4.8 research program is separated
in Section~\ref{v480:integration}. Neither section promotes a finite
detector result to exact all-readout closure.

\paragraph{Current boundary.}
The baseline permits guarded local/restricted and effective
matter-detector statements under their stated assumptions. Exact
finite C5 zero, unrestricted all-readout terminality, C6/T2 closure
and a global photon-clock/Riemann theorem are not established here.
The inherited scale and zero-status distinctions are retained in
Appendix~\ref{v480:support-frames}; the next metric/field-equation
program does not remove them.

\paragraph{Formal-language route.}
QAM supplies object coding and readout records, MML an instruction
language, and MMA guarded transitions. Their full inherited definitions
are in Appendix~\ref{v480:support-qam}. The later shifted normal form,
not an unqualified return to the pre-shift syntax, is the integration
interface for new material; see
Appendix~\ref{v480:support-shift}. Historical pre-shift displays remain
available with their occurrence and backtranslation records.
Typed commands and short references are neither background executions
nor proof certificates.

\paragraph{Release discipline.}
The published v4.7.0 and v4.8.0 TeX and PDF are unchanged. The current
public baseline remains v4.8.0 (DOI 10.5281/zenodo.23365244).
No v4.9.0 version DOI, public upload or new physical closure claim is
assigned to this research candidate. In inherited
fragments, the macro \texttt{docversion} deliberately retains its
v4.7.0 value so historical records are not relabelled as new results.
The working title and metadata identify this copy separately.

\clearpage
